% ############################################################################
% ##  SUPERSEDED 2026-08-03 -- NOT COMPILED BY ANY DOCUMENT. DO NOT EDIT.
% ##
% ##  This file is a per-theorem rendering from an ABANDONED architecture in
% ##  which paper/main.tex would \input one file per result.  paper/main.tex
% ##  DOES NOT EXIST.  The manuscript is paper/qmd/{main,supplement}.qmd, and
% ##  both \input the theory corpus directly.  Verified 2026-08-03: neither
% ##  .qmd references sections/theorems.
% ##
% ##  The text below is ROUND-2 vintage against a corpus now at dispatch 25.
% ##  66-78% of its substantive lines no longer match the corpus -- that is old
% ##  text the corpus has since revised, NOT unique work.  Every \label here is
% ##  also declared in the live corpus.
% ##
% ##  SOURCE OF RECORD:
% ##    statements  quality_reports/theory/panel-npiv-discrete/results.tex
% ##    assumptions quality_reports/theory/panel-npiv-discrete/assumptions.tex
% ##    proofs      quality_reports/theory/panel-npiv-discrete/proofs.tex
% ##
% ##  Kept rather than deleted because nothing in this repo is committed.
% ##  Full rationale: paper/sections/theorems/_SUPERSEDED.md
% ############################################################################

% ============================================================================
%  lem:bierens -- equivalence of the Delta-Z-conditional restriction (A) with
%  the unconditional Bierens family {M(u,t)}.
%  Project: panel-npiv-discrete.  Target of plan Sec.2.7; memo Sec.10 Step 2.
%  ROUND 2: revised against the dispatch-1 critic report (56/100).  THAT REPORT
%  IS NO LONGER AT THE PATH THIS HEADER USED TO NAME.  It was written as
%  "proofs_review.md"; that file was later split into
%    quality_reports/theory/panel-npiv-discrete/proofs_review_dispatch1.md
%    quality_reports/theory/panel-npiv-discrete/proofs_review_dispatch2.md
%  and no file called proofs_review.md exists.  A header that cites a dead path
%  is the hand-maintained-header failure mode again (see the SYNC DISCIPLINE
%  note in lem-adjoint.tex): it looks verifiable and is not.
%
%  ROUND-7 SYNC, 2026-07-29 -- THIS FILE IS RESYNCED ON ONE POINT AND ONE ONLY.
%  ------------------------------------------------------------------------
%  psi -> \mathsf L_\Gamma in the proof of lem:bierenscore (Steps 1-4).  Reason:
%  psi is RESERVED for the moment function
%      psi_i(theta, varphi, L) = Delta Z_i Delta U_i + Wtilde_i L(Delta U_i, Delta Z_i)
%  (notation glossary Sec.4, C26), and rem-nophifree.tex uses it in exactly that
%  reserved sense at seven sites.  SINCE ALL FOUR THEOREM FILES NOW COMPILE INTO
%  ONE DOCUMENT, leaving the Laplace transform as psi put two meanings of one
%  glyph inside a single artifact -- a live INV-7 collision, not a latent one.
%  L_Gamma names the two-sided Laplace transform of the image measure Gamma
%  built in Step 1, so NO NEW LETTER is introduced; sans-serif matches the
%  \mathsf c / \mathsf k / \mathsf r family (glossary Sec.1.6).
%
%  NOTHING ELSE IN THIS FILE IS RESYNCED.  The statements and proofs below are
%  the round-2 text and are correct as such; the Part C repairs of rounds 3-7
%  (if any) are not reflected.  When Parts C, E and F are generated together
%  after Step 4, this file must be regenerated MECHANICALLY from results.tex and
%  proofs.tex and diffed, not retyped.
%
%  SCOPE.  (a) The F4 version only: per plan Sec.19.3 the strict-inclusion half
%  -- R(A*_cond) strictly containing the range of the unconditional adjoint --
%  is DROPPED and is not claimed here (see rem:Tcond).  (b) The weight is
%  w = exp, per B1(ii).  A general real-analytic non-polynomial weight is NOT
%  admissible against a bounded ODD transform; see rem:swintercept for the
%  counterexample and for the intercept-indexed family that repairs it.
%
%  Source of record: quality_reports/theory/panel-npiv-discrete/{results,proofs}.tex
%  Assumptions cited by tag; see .../assumptions.tex.
%  Requires \IR, \un, \upmodels, \vdagger from preamble.tex.
%
%  DEPENDENCIES that must appear EARLIER in the compiled document:
%    def:operator   -- M(u,t;varphi,theta) and T(varphi,theta), and Pi;
%    def:admissible -- I and V_CI;
%    lem:quantu     -- L-3, the binary identity and the single-null-set
%                      extension (cited in Steps 3 and 5 below);
%    lem:density    -- L-2 (cited in rem:Tcond only);
%    lem:shiftinv   -- cited in the proof of cor:opsol.
%  All five are in results.tex.  This file is not self-contained by design:
%  the supporting lemmas belong to the section body, not to the theorem file.
%
%  THIS FILE IS AN ALTERNATIVE RENDERING of lem:bierenscore, lem:bierens and
%  cor:opsol, carrying the same \label keys as results.tex.  The two must never
%  be \input into the same document.
% ============================================================================

\begin{lemma}\label{lem:bierenscore}
A bounded, one-to-one transform and the exponential weight convert a conditional moment restriction into
an unconditional family indexed by a set of positive Lebesgue measure.
Suppose Assumption B1(i)--(iii) holds, $V$ is a random vector in $\IR^{p}$, and $\varepsilon$ is a random
variable with $E|\varepsilon|<\infty$. Then
\[
E\big[\varepsilon\mid V\big]=0\ \text{a.s.}
\qquad\Longleftrightarrow\qquad
E\big[\varepsilon\,e^{t'\Phi(V)}\big]=0\ \ \text{for every }t\in\mathcal T .
\]
\end{lemma}

\begin{proof}
The forward implication is immediate; the converse is proved in five steps.

\emph{Step 0: the forward implication.}
Suppose $E[\varepsilon\mid V]=0$ a.s. By B1(i)--(ii), $|e^{t'\Phi(V)}|\le\bar w(t)<\infty$ for every $t$,
so $E|\varepsilon\,e^{t'\Phi(V)}|\le\bar w(t)E|\varepsilon|<\infty$ and the tower property gives
$E[\varepsilon\,e^{t'\Phi(V)}]=E\big[E[\varepsilon\mid V]\,e^{t'\Phi(V)}\big]=0$. Boundedness of $\Phi$
is doing the work: without it the product need not be integrable and the step fails absent a moment
condition on $V$.

\emph{Step 1: a compactly supported finite signed measure.}
Define $\Gamma(B):=E[\varepsilon\,\un\{\Phi(V)\in B\}]$ for Borel $B\subseteq\IR^{p}$. Then $\Gamma$ is a
finite signed measure with $|\Gamma|(\IR^{p})\le E|\varepsilon|<\infty$, and by B1(i) it is supported in
the compact cube $[-c_\Phi,c_\Phi]^{p}$. By the change-of-variables formula for image measures,
\begin{equation}\label{eq:laplace}
\mathsf L_\Gamma(t)\;:=\;E\big[\varepsilon\,e^{t'\Phi(V)}\big]\;=\;\int_{[-c_\Phi,c_\Phi]^{p}}e^{t'v}\,\Gamma(dv),
\qquad t\in\IR^{p},
\end{equation}
absolutely convergent for every $t$ because $|e^{t'v}|\le e^{c_\Phi\|t\|_1}$ on the cube.

\emph{Step 2: $\mathsf L_\Gamma$ is real-analytic on $\IR^{p}$.}
Write $\mathbf j=(j_1,\dots,j_p)$ for a tuple of nonnegative integers, with
$t^{\mathbf j}:=\prod_{q\le p}t_q^{j_q}$, $|\mathbf j|:=\sum_q j_q$, $\mathbf j!:=\prod_q j_q!$ and
$\partial^{\mathbf j}:=\partial^{j_1}_{t_1}\cdots\partial^{j_p}_{t_p}$ --- \emph{the multi-index carries
no symbol, $\alpha$ being binding as the propensity rate exponent (glossary C72)}. Expand
$e^{t'v}=\sum_{\mathbf j}t^{\mathbf j}v^{\mathbf j}/\mathbf j!$. Since
$|v^{\mathbf j}|\le c_\Phi^{|\mathbf j|}$ on the cube,
$\sum_{\mathbf j}(|t^{\mathbf j}|/\mathbf j!)\int|v^{\mathbf j}|\,|\Gamma|(dv)\le|\Gamma|(\IR^{p})e^{c_\Phi\|t\|_1}<\infty$
for every $t$, so the series converges absolutely everywhere and may be integrated term by term in
\eqref{eq:laplace} by dominated convergence. Hence $\mathsf L_\Gamma$ is real-analytic on $\IR^{p}$ with
$\partial^{\mathbf j}\mathsf L_\Gamma(0)=\int v^{\mathbf j}\,\Gamma(dv)$.

\emph{Step 3: $\mathsf L_\Gamma$ vanishes identically.}
By hypothesis $\mathsf L_\Gamma=0$ on $\mathcal T$, and $\mathrm{Leb}(\mathcal T)>0$ by B1(iii). A
real-analytic function on $\IR^{p}$ --- which is connected and open --- whose zero set has strictly
positive Lebesgue measure vanishes identically on $\IR^{p}$; this is the identity theorem for
real-analytic functions and is the step made in \citet[proof of Thm.~1]{bierens1990}. Hence
$\mathsf L_\Gamma\equiv0$ on $\IR^{p}$.

\emph{Step 4: all moments of $\Gamma$ vanish, hence $\Gamma=0$.}
By Steps 2 and 3, $\int v^{\mathbf j}\,\Gamma(dv)=0$ for every $\mathbf j$, so $\int P\,d\Gamma=0$ for every
polynomial $P$. By the Stone--Weierstrass theorem polynomials are uniformly dense in
$C([-c_\Phi,c_\Phi]^{p})$ and $|\Gamma|$ is finite, so $\int f\,d\Gamma=0$ for every continuous $f$ on
the cube; by the Riesz representation theorem for finite signed Borel measures on a compact metric space,
$\Gamma=0$. This is the second place where boundedness of $\Phi$ is load-bearing: it makes $\Gamma$
compactly supported, hence determined by its moments.

\emph{Step 5: back to $V$.}
$\Gamma=0$ means $E[\varepsilon\mid\Phi(V)]=0$ a.s. By B1(i), $\Phi$ is an injective Borel map between
Polish spaces, hence bi-measurable by the Lusin--Souslin theorem, so $\sigma(\Phi(V))=\sigma(V)$ and
$E[\varepsilon\mid V]=0$ a.s. (For the maintained choice $\Phi$ is a homeomorphism onto an open cube and
bi-measurability is elementary.)
\end{proof}

\begin{lemma}\label{lem:bierens}
The $\Delta Z$-conditional restriction \textup{(A)} is equivalent to the vanishing of the unconditional
family $\{\mathcal M(u,t)\}$, so no conditional distribution given $\Delta Z$ ever has to be estimated.
Suppose Assumptions B1(i)--(iv) hold. Then the following are equivalent:
\begin{enumerate}
\item[(a)] $\Delta U\upmodels W\mid\Delta Z$;
\item[(b)] $E[\widetilde W\,\varsigma(\Delta U)\mid\Delta Z]=0$ a.s.\ for every bounded measurable $\varsigma$, that is, restriction \textup{(A)};
\item[(c)] $\mathcal M(u,t)=0$ for every $(u,t)\in\IR\times\mathcal T$;
\item[(d)] $T(\vdagger,\theta_\dagger)=0$ in $L^{2}(\Pi)$.
\end{enumerate}
\end{lemma}

\begin{proof}
The strategy is $(a)\Rightarrow(b)\Rightarrow(c)\Rightarrow(a)$, then $(c)\Leftrightarrow(d)$.
Throughout, $|\widetilde W|\le1$ almost surely --- free from $W\in\{0,1\}$ and $\varrho\in[0,1]$, with
no overlap condition used --- and $E[\widetilde W\mid\Delta Z]=0$ by the definition of $\varrho$.

\emph{Step 1: $(a)\Rightarrow(b)$.}
Let $\varsigma$ be bounded and measurable. Under $(a)$, $\varsigma(\Delta U)$ and $W$ are conditionally
independent given $\Delta Z$, so
$E[\widetilde W\varsigma(\Delta U)\mid\Delta Z]=E[\widetilde W\mid\Delta Z]E[\varsigma(\Delta U)\mid\Delta Z]=0$.

\emph{Step 2: $(b)\Rightarrow(c)$.}
Fix $u$ and take $\varsigma=\un\{\cdot\le u\}$ in $(b)$. Fix $t\in\mathcal T$; the product
$\widetilde W\un\{\Delta U\le u\}e^{t'\Phi(\Delta Z)}$ is integrable by B1(i)--(ii), and
$e^{t'\Phi(\Delta Z)}$ is $\sigma(\Delta Z)$-measurable, so the tower property gives
$\mathcal M(u,t)=E\big[e^{t'\Phi(\Delta Z)}\,E[\widetilde W\un\{\Delta U\le u\}\mid\Delta Z]\big]=0$.

\emph{Step 3: $(c)\Rightarrow(a)$.}
Fix $u\in\IR$ and set $\varepsilon:=\widetilde W\un\{\Delta U\le u\}$, so $|\varepsilon|\le1$ and
$E|\varepsilon|<\infty$. By $(c)$, $E[\varepsilon\,e^{t'\Phi(\Delta Z)}]=0$ for every $t\in\mathcal T$,
so \cref{lem:bierenscore} with $V=\Delta Z$ gives $E[\widetilde W\un\{\Delta U\le u\}\mid\Delta Z]=0$
a.s. This holds for every $u\in\mathbb{Q}$, so by \cref{lem:quantu} applied at
$(\theta,\varphi)=(\theta_\dagger,\vdagger)$, and writing $F_j:=F_j^{(\theta_\dagger,\vdagger)}$, there
is a single null set $N$ outside which $\varrho(z)(1-\varrho(z))\{F_1(u\mid z)-F_0(u\mid z)\}=0$ for
every $u\in\IR$. Fix $z\notin N$ and split on
$\mathcal Z_0:=\{z:\varrho(z)(1-\varrho(z))=0\}$. If $z\notin\mathcal Z_0$ the
prefactor is strictly positive, so $F_1(\cdot\mid z)=F_0(\cdot\mid z)$, which for binary $W$ is exactly
conditional independence at $z$. If $z\in\mathcal Z_0$ then $P(W=1\mid\Delta Z=z)\in\{0,1\}$, so $W$ is
conditionally degenerate at $z$ and conditional independence holds there automatically in the
factorisation sense. Combining gives $(a)$. \emph{No overlap condition is used}: on the degenerate region
$\widetilde W=0$ almost surely, so the restriction has no bite there --- but $(a)$ holds there regardless,
so there is nothing to recover.

\emph{Step 4: $(c)\Rightarrow(d)$.} Immediate.

\emph{Step 5: $(d)\Rightarrow(c)$.}
$T(\vdagger,\theta_\dagger)=0$ in $L^{2}(\Pi)$ means $\mathcal M=0$ $\Pi$-a.e. Since
$\Pi=\Pi_u\otimes\Pi_t$, Fubini's theorem supplies a $\Pi_t$-null set $N_{\mathcal T}$ such that for
every $t\notin N_{\mathcal T}$, $\mathcal M(u,t)=0$ for $\Pi_u$-a.e.\ $u$. Fix such a $t$ and let
$Z_t:=\{u\in\IR:\mathcal M(u,t)=0\}$, so $\Pi_u(Z_t^{c})=0$. If $Z_t$ were not dense in $\IR$ there
would be a nonempty open $O\subseteq Z_t^{c}$ with $\Pi_u(O)=0$, contradicting the full topological
support of $\Pi_u$ assumed in B1(iv.a); so $Z_t$ is dense. The map $u\mapsto\mathcal M(u,t)$ is
right-continuous, since $u_k\downarrow u$ gives $\un\{\Delta U\le u_k\}\downarrow\un\{\Delta U\le u\}$
pointwise with integrand dominated by $\bar w(t)$. Choosing $u_k\in Z_t$ with $u_k\downarrow u$ gives
$\mathcal M(\cdot,t)\equiv0$ on $\IR$.

By B1(iv.b) there is a Borel $\mathcal T_0\subseteq\mathcal T$ with $\mathrm{Leb}|_{\mathcal T_0}$ and
$\Pi_t|_{\mathcal T_0}$ not mutually singular. If $\mathrm{Leb}(\mathcal T_0\setminus N_{\mathcal T})=0$
then $\mathrm{Leb}|_{\mathcal T_0}$ is concentrated on $N_{\mathcal T}$ while $\Pi_t|_{\mathcal T_0}$ is
concentrated on its complement, making the two mutually singular --- a contradiction. Hence
$\mathrm{Leb}(\mathcal T_0\setminus N_{\mathcal T})>0$, and running Step 3 with $\mathcal T$ replaced by
$\mathcal T_0\setminus N_{\mathcal T}$ --- legitimate because \cref{lem:bierenscore} requires only
positive Lebesgue measure of the index set --- delivers $(a)$. Then $(c)$ on all of $\mathcal T$ follows
from Steps 1 and 2.
\end{proof}

\begin{corollary}[to \cref{lem:shiftinv} and \cref{lem:bierens}]\label{cor:opsol}
The operator equation has the same solution set in $\varphi$ at every $\theta$, and that set is
$\mathcal V_{CI}$.
Suppose Assumptions B1(i)--(iv) hold. Then for every $\theta\in\IR^{p}$,
\[
\big\{\varphi\in\mathcal V:\ T(\varphi,\theta)=0\ \text{in }L^{2}(\Pi)\big\}\;=\;\mathcal V_{CI},
\]
and under Assumption A1b(iii) this set equals $\{\vdagger\}$.
\end{corollary}

\begin{proof}
Fix $\theta\in\IR^{p}$ and $\varphi\in\mathcal V$. The proof of \cref{lem:bierens} uses no property of
$(\vdagger,\theta_\dagger)$ beyond three facts: $|\widetilde W|\le1$; $E[\widetilde W\mid\Delta Z]=0$;
and the applicability of \cref{lem:quantu} to the residual. All three hold at arbitrary trial values, and
\cref{lem:quantu} is stated at arbitrary trial values for exactly this reason. Substituting
$\Delta U(\theta,\varphi)$ for $\Delta U$ throughout that proof gives
$T(\varphi,\theta)=0\Leftrightarrow\Delta U(\theta,\varphi)\upmodels W\mid\Delta Z\Leftrightarrow(\varphi,\theta)\in\mathcal I$.
By \eqref{eq:shiftinv} of \cref{lem:shiftinv}, which holds with no assumption,
$(\varphi,\theta)\in\mathcal I$ if and only if $\varphi\in\mathcal V_{CI}$, whatever the value of
$\theta$. Under A1b(iii) that set is $\{\vdagger\}$.
\end{proof}

% NEITHER A1 NOR A4(i) IS A HYPOTHESIS of lem:bierens or cor:opsol.
% A1 would make lem:bierens VACUOUS: A1's first clause IS statement (a), so
% under A1 all four of (a)-(d) are true and the printed equivalence asserts
% nothing.  A4(i) is idle: see Step 3.  What this buys is worth a sentence in
% the paper -- the equivalence needs NO MOMENT CONDITION on Delta Y, Delta U or
% Delta Z, and NO LEBESGUE DENSITY for Delta U or Delta Z.  That is what
% boundedness of Phi purchases; see rem:Tcond.

\begin{remark}\label{rem:swintercept}
B1(ii) fixes $w=\exp$, and the restriction is substantive rather than expositional.

\emph{The textbook condition is not enough.} Suppose $w$ were required only to be real-analytic and
non-polynomial, as in the usual statement of the device. Take $V=\pm1$ with probability $\tfrac12$ each,
let $\Phi$ be the maintained componentwise $\arctan$ --- an \emph{odd} function, so
$\Phi(-1)=-\Phi(1)$
--- and set $\varepsilon:=V$. Then $E|\varepsilon|<\infty$ and $E[\varepsilon\mid V]=V\neq0$, yet for
$w=\cos$ and \emph{every} $t$,
\[
E\big[\varepsilon\,\cos\big(t\Phi(V)\big)\big]
\;=\;\tfrac12\cos\big(t\,\Phi(1)\big)-\tfrac12\cos\big(-t\,\Phi(1)\big)\;=\;0 .
\]
B1(i)--(iii) hold and the converse direction of \cref{lem:bierenscore} fails. The mechanism is that an
odd $\Phi$ against an even $w$ exposes only the even part of the image measure, so Step 4 of the proof
recovers the even moments and no others. The failure is not confined to the abstract lemma: with
$\Delta Z$ scalar and two-valued at $\pm c$,
$\mathcal M(u,t)=\cos\big(t\,\Phi(c)\big)\,E[\widetilde W\un\{\Delta U\le u\}]$, so (c) of
\cref{lem:bierens} holds
whenever the \emph{unconditional} restriction does, while the $\Delta Z$-conditional restriction can fail
through opposite-signed violations that cancel across the two values of $\Delta Z$.

\emph{The correct general statement, and what it would cost.} The repair is not to strengthen $\Phi$ but
to add an intercept. Index the family by $(t_0,t)$ over a subset of $\IR^{1+p}$ of positive Lebesgue
measure and use $\{w(t_0+t'\Phi(z))\}$. For $w$ real-analytic and non-polynomial, each derivative
$w^{(k)}$ is analytic and not identically zero --- otherwise $w$ is a polynomial --- so its zero set is
discrete and hence countable; therefore all but countably many $c\in\IR$ satisfy $w^{(k)}(c)\neq0$ for
every $k\ge0$, and the argument runs at such a $c$. This is the Stinchcombe--White family and the result
is \citet[Thm.~3.1]{stinchcombewhite1998}. Adopting it would change the range space to
$L^{2}(\IR\times\mathcal T^{+},\Pi)$ with $\mathcal T^{+}\subseteq\IR^{1+p}$, and would require the
envelope bookkeeping of B1(iv)--(v) to be redone with
$\bar w(t_0,t)=\sup_{|v|\le|t_0|+c_\Phi\|t\|_1}|w(v)|$.

\emph{Why it is not adopted.} $w=\exp$ is the maintained choice of the estimator, it is what
\citet[Thm.~1]{bierens1990} characterises, and the additional generality is consumed nowhere in the
programme.
\end{remark}

\begin{remark}\label{rem:Tcond}
Compactness of $\mathcal T$ is \emph{sufficient but not necessary} and should not be assumed. What the
proofs use is B1(iv.c) and B1(v), which for $w=\exp$ read
$\Pi_u(\IR)\int e^{2c_\Phi\|t\|_1}\Pi_t(dt)<\infty$ and $\int e^{c_\Phi\|t\|_1}|\gamma|(dt)<\infty$; a
compact $\mathcal T$ carrying a finite $\Pi_t$ and a finite $\gamma$ satisfies both, but so do many
unbounded index sets. Compactness was introduced in order to license differentiation under the
$t$-integral, and B1(v) supplies that directly: differentiating
$L(u,z)=\int_{\mathcal T}\lambda(u,t)w(t'\Phi(z))\gamma(dt)$ in $u$ under the integral sign is licensed
by $|\partial_u\lambda(u,t)\,w(t'\Phi(z))|\le\Lambda_1(t)\bar w(t)$ together with
$\int_{\mathcal T}\Lambda_1\bar w\,d|\gamma|<\infty$.

Two clauses of B1(iv) are minimal and should not be strengthened back. Step 5 above consumes only
(iv.a), full topological support of $\Pi_u$, and (iv.b), non-mutual-singularity of
$\mathrm{Leb}|_{\mathcal T_0}$ and $\Pi_t|_{\mathcal T_0}$ --- exactly the statement that no
$\Pi_t$-null set can swallow $\mathcal T_0$. By the Lebesgue decomposition, (iv.b) is equivalent to
$\Pi_t$ having a nonzero absolutely continuous part, so it excludes \emph{every} Lebesgue-singular
$\Pi_t$ and not merely the atomic ones: a Cantor measure is non-atomic and is excluded. It is also
sufficient rather than necessary. What the argument needs is an analytic-uniqueness condition on
$\operatorname{supp}\Pi_t$: by Step 2 of the proof of \cref{lem:bierenscore}, $t\mapsto\mathcal M(u,t)$ is
real-analytic on $\IR^{p}$, so it suffices that its zero set cannot contain $\operatorname{supp}\Pi_t$ up to
a $\Pi_t$-null set unless it vanishes identically. For $p=1$ this holds whenever the support has an
accumulation point --- a Cantor measure included, since a Cantor set is perfect --- and
\cref{lem:bierens} survives such a $\Pi_t$; for $p\ge2$ the requirement is that the support not lie in a
proper real-analytic variety, which perfectness alone does not give. What (iv.b) excludes in every
dimension is a finite grid. \emph{That was the ground on which the growing-family condition B1(vi) was
classed an identification assumption rather than an approximation device, and B1(vi) is now
\textbf{retired}} --- not because the classification was wrong but because the clause describes an estimator
this programme does not implement: the operator integrates against $\Pi_t$ over all of $\mathcal T$ and never
over a finite subfamily. \emph{Identification of $\vdagger$ here is A1b(iii), assumed outright}, and the
sentence is kept in the corrected register rather than deleted so that the exclusion (iv.b) performs stays
visible to a reader who reinstates the growing-grid design.

It is also worth recording what the equivalence does \emph{not} need: no moment condition on $\Delta Y$,
$\Delta U$ or $\Delta Z$, and no Lebesgue density for $\Delta U$ or $\Delta Z$. That is what boundedness
of $\Phi$ purchases.

The correction measure is in any case nearly moot on the $\theta$-side: by \cref{lem:density} the
estimator \emph{solves for} $L$ inside a dense class rather than integrating a fixed weighting measure
against a continuum, so no $t$-integral survives in the construction of $\hat L$. This is also the
surviving argument against a continuum-GMM formulation: that would require choosing a weighting measure
on $\mathcal T$ and regularising the inverse of its covariance operator, and the present construction
requires no weighting-measure choice at all.

Finally, \cref{lem:bierens} is stated at the equivalence, not at a strict inclusion. The claim that
$\mathcal R(\mathcal A^\ast_{\mathrm{cond}})$ strictly contains the range of the unconditional adjoint is
not proved here and is not claimed: with $W$ binary the two moment classes are not nested, since
$\{g:E[g\mid\Delta U,\Delta Z]=0\}=\{\widetilde W\ell(\Delta U,\Delta Z)\}$ is an equality of sets, and
the comparison across the two modelling assumptions is heuristic. The conditional route is defended on
economic grounds and on internal consistency with the maintained exogeneity assumption, not on a weaker
source condition.
\end{remark}
